% ============================================================
% Gesamtskript VWL --- Musterlösungen Lektion 3
% Konjunktur, Wirtschaftspolitik & Außenwirtschaft (Rahmenplan DIHK 1.1)
% Punktangaben in eckigen Klammern = Bewertungsvorschlag
% ============================================================

\begin{loesung}[title={Lösung Aufgabe~3.1 --- Konjunkturlage richtig deuten}]
\textbf{a) Zuordnung} \hfill [5~P.: je richtige Zuordnung 1~P.]

\begin{tabular}{p{7.6cm} p{5.6cm}}
\toprule
\textbf{Datum} & \textbf{Indikatorart}\\
\midrule
(1) ifo-Geschäftsklimaindex & Frühindikator\\
(2) Auftragseingänge & Frühindikator\\
(3) Kapazitätsauslastung & Präsenzindikator\\
(4) Arbeitslosenquote & Spätindikator\\
(5) Unternehmensinsolvenzen & Spätindikator\\
\bottomrule
\end{tabular}

\medskip
\textbf{b) Konjunkturphase} \hfill [3~P.]

Die Volkswirtschaft befindet sich am \emph{Ende der Tiefphase (Depression/Talsohle)} bzw.
am \emph{Übergang zum Aufschwung} (unterer Wendepunkt). [1~P.]

Begründung: Die \emph{Präsenz}- und \emph{Spätindikatoren} zeigen noch eine schwache Lage
(unterdurchschnittliche Auslastung, steigende Arbeitslosigkeit, viele Insolvenzen). Die
\emph{Frühindikatoren} (Geschäftsklima, Auftragseingänge) steigen jedoch bereits und
kündigen eine Belebung an. Da Spätindikatoren der Entwicklung hinterherlaufen, ist ihr
negativer Verlauf mit einer beginnenden Erholung vereinbar. [2~P.]

\medskip
\textbf{c) Schlussfolgerungen} \hfill [3~P.: zwei Schlussfolgerungen, schlüssig erläutert]

Erwartet werden z.\,B. (zwei genügen):
\begin{itemize}
  \item \emph{Personal:} Fachkräfte halten statt entlassen (ggf. Kurzarbeit nutzen und
    schrittweise zurückfahren), Einstellungen für steigende Auftragszahlen vorbereiten
    --- im Aufschwung wird der Arbeitsmarkt schnell enger.
  \item \emph{Produktion/Kapazität:} Kapazitäten so planen, dass steigende Aufträge
    bedient werden können (freie Kapazitäten zunächst auslasten, Erweiterungen vorbereiten).
  \item \emph{Investitionen/Finanzierung:} Investitionen jetzt planen, solange Zinsen und
    Preise für Investitionsgüter in der Tiefphase noch vergleichsweise niedrig sind.
  \item \emph{Einkauf/Lager:} Material frühzeitig disponieren, da im Aufschwung Preise
    und Lieferzeiten steigen.
\end{itemize}
\end{loesung}

\begin{loesung}[title={Lösung Aufgabe~3.2 --- Inflation, Realzins und Lohnabschluss}]
\textbf{a) Inflationsrate} \hfill [3~P.]
\[
\text{Inflationsrate} = \frac{127{,}1 - 124{,}0}{124{,}0} \times 100
= \frac{3{,}1}{124{,}0} \times 100 = 0{,}025 \times 100 = \textbf{2{,}5~\%}
\]
Probe: $124{,}0 \times 1{,}025 = 127{,}1$. \emph{Typischer Fehler:} Punktdifferenz
(3,1~Punkte) mit der Inflationsrate gleichsetzen --- es muss durch den alten Indexstand
geteilt werden.

\medskip
\textbf{b) Zinsertrag und Realzins} \hfill [4~P.]

Zinsertrag: $250.000~\text{€} \times 0{,}034 = \textbf{8.500~€}$ [1~P.]

Realzins (Näherung): Nominalzins $-$ Inflationsrate $= 3{,}4~\% - 2{,}5~\% =
\textbf{0{,}9~\%}$ [2~P.]\\
{\small (Exakt: $1{,}034 \div 1{,}025 - 1 \approx 0{,}0088 = 0{,}88~\%$ --- beide Wege
werden anerkannt.)}

Beurteilung: Der reale Vermögenszuwachs ist gering. Von den 8.500~€ Zinsen gleicht der
größte Teil nur den Kaufkraftverlust aus ($250.000 \times 0{,}025 = 6.250$~€); nur rund
0,9~\% sind echter Kaufkraftgewinn (vor Steuern). Nach Abzug von Steuern kann der
Realzins sogar negativ werden. [1~P.]

\medskip
\textbf{c) Nominal- und Reallohn} \hfill [3~P.]

Nominale Steigerung: $\dfrac{3.264 - 3.200}{3.200} \times 100 = \dfrac{64}{3.200} \times 100
= \textbf{2{,}0~\%}$ [1~P.]

Reallohnänderung (Näherung): $2{,}0~\% - 2{,}5~\% = \textbf{$-$0{,}5~\%}$ [1~P.]\\
{\small (Exakt: $1{,}020 \div 1{,}025 - 1 \approx -0{,}49~\%$.)}

Interpretation: Obwohl der Disponent nominal mehr verdient, kann er sich real
\emph{weniger} kaufen, weil die Preise stärker gestiegen sind als sein Gehalt
(Kaufkraftverlust). [1~P.]

\medskip
\textbf{d) Verteilungsspielraum und Beurteilung} \hfill [4~P.]

Verteilungsspielraum $=$ Produktivitätsfortschritt $+$ Zielinflation
$= 1{,}3~\% + 2{,}0~\% = \textbf{3{,}3~\%}$ [1~P.]\\
Die Forderung liegt $5{,}5 - 3{,}3 = \textbf{2{,}2~Prozentpunkte}$ über dem Spielraum.
[1~P.]

\emph{Sicht der Arbeitgeber} (produktivitätsorientierte Lohnpolitik): Ein Abschluss über
3,3~\% erhöht die Lohnstückkosten. Folgen: Preiserhöhungen (Lohn-Preis-Spirale,
Inflationsdruck), sinkende Wettbewerbsfähigkeit, Rationalisierung oder Verlagerung
--- die Forderung ist zu hoch. [1~P.]

\emph{Sicht der Gewerkschaft} (Kaufkrafttheorie): Die tatsächliche Inflation (2,5~\%)
liegt über dem Ziel; die Beschäftigten haben real verloren (vgl.~c). Höhere Löhne
sichern die Kaufkraft, stärken die Nachfrage und damit Produktion und Beschäftigung
--- die Forderung ist berechtigt. [1~P.]
\end{loesung}

\begin{loesung}[title={Lösung Aufgabe~3.3 --- Wechselkursänderung im Exportgeschäft}]
\emph{Rechenregel Mengennotierung (1~€ $=$ x~USD):} Europreis $\times$ Kurs $=$
Dollarpreis; Dollarpreis $\div$ Kurs $=$ Europreis.

\medskip
\textbf{a) Umrechnung zum Kurs 1,08} \hfill [2~P.: je 1~P.]

Messgerät: $45.000~\text{€} \times 1{,}08 = \textbf{48.600~USD}$\\
Sensoren: $32.400~\text{USD} \div 1{,}08 = \textbf{30.000~€}$

\medskip
\textbf{b) Kurs 1,20} \hfill [5~P.]

Für 1~€ erhält man nun mehr US-Dollar (1,20 statt 1,08) $\Rightarrow$ der Euro hat
\textbf{aufgewertet} (der Dollar abgewertet). [1~P.]

Neuer Dollarpreis: $45.000~\text{€} \times 1{,}20 = \textbf{54.000~USD}$ [1~P.]\\
Veränderung: $\dfrac{54.000 - 48.600}{48.600} \times 100 = \dfrac{5.400}{48.600} \times 100
\approx \textbf{+11,11~\%}$ [1~P.]

Sensoren: $32.400~\text{USD} \div 1{,}20 = \textbf{27.000~€}$, also
$3.000$~€ bzw. $\dfrac{3.000}{30.000} \times 100 = 10~\%$ günstiger. [2~P.]

Fazit: Die Aufwertung verteuert die Exporte für US-Kunden und verbilligt die Importe
für Nordlicht Optik.

\medskip
\textbf{c) Unveränderter Dollarpreis} \hfill [2~P.]

Euro-Erlös: $48.600~\text{USD} \div 1{,}20 = \textbf{40.500~€}$ [1~P.]\\
Erlösminderung: $45.000 - 40.500 = \textbf{4.500~€}$, das sind
$\dfrac{4.500}{45.000} \times 100 = \textbf{10~\%}$ je Gerät. [1~P.]

\medskip
\textbf{d) Absicherung} \hfill [2~P.: je Nennung 1~P.]

Z.\,B.: Devisentermingeschäft (Kurs wird heute für einen späteren Zahlungstermin
festgeschrieben); Devisenoption; Fakturierung in Euro (Risiko trägt der Kunde);
Kurssicherungsklauseln im Vertrag; „natürliche Absicherung“, indem Einkäufe (z.\,B.
Sensoren) und Verkäufe in derselben Währung abgewickelt werden.

\medskip
\textbf{e) Lieferungen nach Österreich} \hfill [3~P.]

Österreich gehört zur Eurozone; Lieferungen werden in derselben Währung (Euro) abgerechnet
--- es gibt keinen Wechselkurs und damit weder Kursrisiko noch Umtauschkosten. [1~P.]

Weitere Vorteile (zwei genügen, je 1~P.): keine Zölle und keine mengenmäßigen
Beschränkungen im Binnenmarkt (freier Warenverkehr); harmonisierte technische Normen
(z.\,B. CE-Kennzeichnung), dadurch keine zusätzlichen Zulassungen; Preistransparenz und
einfache Kalkulation; schneller, günstiger Zahlungsverkehr (SEPA); Monteure können
ohne Grenzkontrollen und Arbeitserlaubnis eingesetzt werden (freier Personen- und
Dienstleistungsverkehr).
\end{loesung}

\begin{loesung}[title={Lösung Aufgabe~3.4 --- Regionaler Arbeitsmarkt}]
\textbf{a) Arbeitslosenquoten} \hfill [4~P.]

Formel: Arbeitslosenquote $= \dfrac{\text{registrierte Arbeitslose}}{\text{zivile
Erwerbspersonen}} \times 100$ mit Erwerbspersonen $=$ Erwerbstätige $+$ Arbeitslose.

Vorjahr: Erwerbspersonen $= 236.800 + 13.200 = 250.000$;\quad
$\dfrac{13.200}{250.000} \times 100 = \textbf{5,28~\%}$ [1,5~P.]

Berichtsjahr: Erwerbspersonen $= 234.000 + 16.000 = 250.000$;\quad
$\dfrac{16.000}{250.000} \times 100 = \textbf{6,40~\%}$ [1,5~P.]

Veränderung: $6{,}40 - 5{,}28 = \textbf{+1,12~Prozentpunkte}$ [1~P.]\\
{\small \emph{Achtung:} Prozentpunkte, nicht Prozent! (Relativ wäre der Anstieg der
Quote $1{,}12 \div 5{,}28 \approx 21{,}2~\%$.) Häufiger Fehler: nur durch die
Erwerbstätigen statt durch die Erwerbspersonen teilen.}

\medskip
\textbf{b) Formen der Arbeitslosigkeit} \hfill [4~P.: je Fall 1~P.]

\begin{tabular}{p{1.1cm} p{2.6cm} p{9.4cm}}
\toprule
\textbf{Fall} & \textbf{Form} & \textbf{Begründung}\\
\midrule
(1) & saisonal & jahreszeitlich bedingt (Witterung im Baugewerbe), kehrt regelmäßig wieder\\
(2) & konjunkturell & Ursache ist ein allgemeiner Nachfragerückgang (Rezession)\\
(3) & strukturell & dauerhafter Wandel (technischer Fortschritt); Qualifikation passt nicht mehr zur Nachfrage\\
(4) & friktionell & kurze Such- bzw. Übergangsphase beim Arbeitsplatzwechsel\\
\bottomrule
\end{tabular}

\medskip
\textbf{c) Anteil struktureller Arbeitslosigkeit} \hfill [2~P.]

$\dfrac{7.400}{16.000} \times 100 = \textbf{46,25~\%}$ [1~P.]

Fast die Hälfte der Arbeitslosigkeit ist strukturell bedingt. Sie verschwindet auch im
Aufschwung nicht von selbst, weil Qualifikation, Branche oder Region nicht zu den offenen
Stellen passen (Sockelarbeitslosigkeit); sie ist nur langfristig durch Qualifizierung und
Strukturwandel abzubauen. Es droht Langzeitarbeitslosigkeit mit Kaufkraft- und
Steuerausfällen für die Region. [1~P.]

\medskip
\textbf{d) Maßnahmen} \hfill [3~P.]

\emph{Strukturell:} Umschulung bzw. Weiterbildung (z.\,B. zum Mediengestalter),
Mobilitätshilfen, Ansiedlung neuer Branchen durch Strukturpolitik $\Rightarrow$
\textbf{aktive} Arbeitsmarktpolitik. [1,5~P.]

\emph{Konjunkturell:} Kurzarbeitergeld (Arbeitsplätze werden über die Krise gehalten)
$\Rightarrow$ \textbf{passive} Arbeitsmarktpolitik; ergänzend wirken expansive Geld- und
Fiskalpolitik. [1,5~P.]\\
{\small (Andere sachgerechte Maßnahmen mit korrekter Zuordnung werden anerkannt.)}
\end{loesung}

\begin{loesung}[title={Lösung Aufgabe~3.5 --- Leitzinssenkung der EZB}]
\textbf{a) Begründung der Zinssenkung} \hfill [3~P.]

Vorrangiges Ziel der EZB ist Preisniveaustabilität, definiert als Anstieg des HVPI von
\textbf{2~\% mittelfristig}. Das Ziel ist \textbf{symmetrisch}: Ein dauerhaftes
\emph{Unterschreiten} ist ebenso unerwünscht wie ein Überschreiten. [1~P.]

Bei 1,3~\% und Prognosen deutlich unter 2~\% verfehlt die EZB ihr Ziel nach unten. Es
droht ein Abgleiten in sehr niedrige Inflation oder Deflation: Verbraucher und
Unternehmen schieben Käufe und Investitionen auf, Nachfrage und Beschäftigung sinken
(Abwärtsspirale). [1~P.]

Zudem stagniert die Wirtschaft --- eine expansive Geldpolitik soll die Nachfrage beleben
und die Inflation mittelfristig zum Zielwert zurückführen. [1~P.]

\medskip
\textbf{b) Wirkungskette (expansive Geldpolitik)} \hfill [5~P.: je Stufe 1~P.]

Leitzinssenkung $\Rightarrow$
(1)~Refinanzierung der Geschäftsbanken bei der EZB wird billiger $\Rightarrow$
(2)~Banken geben niedrigere Zinsen weiter, Bauzinsen/Kreditzinsen sinken (Sparen wird
unattraktiver) $\Rightarrow$
(3)~Baufinanzierungen werden günstiger, monatliche Kreditraten sinken; mehr Haushalte
können sich Wohneigentum leisten $\Rightarrow$
(4)~Kreditnachfrage und Nachfrage nach Eigentumswohnungen steigen (auch Kapitalanleger
investieren eher in Immobilien) $\Rightarrow$
(5)~Auftragslage und Umsatz der Wohnbau Lindner GmbH verbessern sich; zugleich sinken
ihre eigenen Finanzierungskosten für Bauprojekte $\Rightarrow$ gesamtwirtschaftlich
steigen Produktion und Beschäftigung, die Inflationsrate nähert sich wieder 2~\%.

\medskip
\textbf{c) Stellungnahme} \hfill [2~P.]

Die Aussage ist falsch. Die Leitzinsen legt der \textbf{EZB-Rat} fest; die EZB ist
\textbf{unabhängig} und an Weisungen der Regierungen nicht gebunden (Geldpolitik). Die
Bundesregierung betreibt Fiskalpolitik (Steuern, Staatsausgaben). [1~P.]\\
Ziel der EZB ist auch nicht die Rettung einzelner Branchen, sondern die
Preisniveaustabilität im gesamten Euroraum; die Bauwirtschaft profitiert nur als
Nebenwirkung. [1~P.]

\medskip
\textbf{d) Grenzen der Geldpolitik} \hfill [2~P.: je Grund 1~P.]

Z.\,B.: Banken geben die Zinssenkung nicht oder nur verzögert weiter (z.\,B. wegen
höherer Risikoaufschläge); pessimistische Erwartungen (Angst vor Arbeitsplatzverlust,
unsichere Konjunktur) --- selbst günstige Kredite werden nicht nachgefragt; hohe Bau- und
Grundstückspreise belasten die Erschwinglichkeit stärker, als die kleine Zinssenkung
entlastet; Wirkungsverzögerungen (Timelags) der Geldpolitik.
\end{loesung}

\begin{loesung}[title={Lösung Aufgabe~3.6 --- Konjunkturprogramm und Schuldenbremse}]
\textbf{a) Zuordnung} \hfill [4~P.: je Maßnahme 0,5~P. Zuordnung, 2~P. Gesamturteil]

\begin{tabular}{p{6.4cm} p{2.4cm} p{4.4cm}}
\toprule
\textbf{Maßnahme} & \textbf{Wirkung} & \textbf{in der Rezession geeignet?}\\
\midrule
(1) Investitionen vorziehen & expansiv & ja (Staatsnachfrage steigt)\\
(2) Umsatzsteuer erhöhen & restriktiv & nein (Kaufkraft sinkt)\\
(3) Sonderabschreibung & expansiv & ja (Steuerentlastung, Investitionsanreiz)\\
(4) Subventionen kürzen, Rücklage & restriktiv & nein (Boom-Maßnahme)\\
\bottomrule
\end{tabular}

\smallskip
Konjunkturgerecht im Sinne der antizyklischen Fiskalpolitik sind in der Rezession nur die
expansiven Maßnahmen \textbf{(1) und (3)}.

\medskip
\textbf{b) Schuldenbremse (Art.~109, 115~GG)} \hfill [4~P.]

Die Schuldenbremse begrenzt die \emph{strukturelle} (konjunkturunabhängige)
Neuverschuldung des Bundes auf 0,35~\% des BIP; die konjunkturbedingte Kreditaufnahme
ist zusätzlich zulässig.

Zulässige strukturelle Neuverschuldung: $4.400~\text{Mrd.~€} \times 0{,}0035 =
\textbf{15,4~Mrd.~€}$ [1~P.]\\
Geplante strukturelle Neuverschuldung: $41 - 22 = \textbf{19~Mrd.~€}$ [1~P.]\\
$\Rightarrow$ $19 - 15{,}4 = \textbf{3,6~Mrd.~€}$ über der Grenze (strukturell
$19 \div 4.400 \times 100 \approx 0{,}43~\%$ des BIP). [1~P.]

Ergebnis: Die Schuldenbremse wird \textbf{nicht eingehalten}. Zulässig wäre die Planung
nur bei einer vom Bundestag festgestellten außergewöhnlichen Notsituation (mit
Tilgungsplan); sonst muss das strukturelle Defizit um 3,6~Mrd.~€ gesenkt werden. [1~P.]

\medskip
\textbf{c) Prozyklische Wirkung und Crowding-out} \hfill [3~P.]

\emph{Prozyklische Wirkung (Timelags):} Zwischen Erkennen der Rezession, Beschluss,
Ausschreibung und Bauausführung vergeht viel Zeit (Erkennungs-, Entscheidungs- und
Wirkungsverzögerung). Setzt die zusätzliche Nachfrage erst ein, wenn die Wirtschaft
schon wieder im Aufschwung ist, verstärkt sie den Boom (Kapazitätsengpässe,
steigende Baupreise) --- die Maßnahme wirkt dann prozyklisch. [2~P.]

\emph{Crowding-out:} Die kreditfinanzierte Staatsnachfrage erhöht die Nachfrage am
Kapitalmarkt und damit die Zinsen; private Investitionen werden dadurch verdrängt, sodass
der gesamtwirtschaftliche Impuls ganz oder teilweise verpufft. [1~P.]
\end{loesung}

\begin{loesung}[title={Lösung Aufgabe~3.7 --- Zahlungsbilanz und Handelskonflikt}]
\textbf{a) Teilbilanzen (in Mrd.~€)} \hfill [4~P.]

Handelsbilanz: $1.480 - 1.290 = \textbf{+190}$ (Überschuss) [1~P.]\\
Dienstleistungsbilanz: $395 - 455 = \textbf{$-$60}$ (Defizit) [1~P.]\\
Leistungsbilanz: $190 + (-60) + 115 + (-72) = \textbf{+173}$ (Überschuss) [2~P.]

\medskip
\textbf{b) Ausgleich der Zahlungsbilanz} \hfill [2~P.]

Jede Transaktion wird doppelt verbucht (Leistung und Gegenleistung). Dem
Leistungsbilanzüberschuss von 173~Mrd.~€ steht daher spiegelbildlich ein
\textbf{Nettokapitalexport} in der \textbf{Kapitalbilanz} in gleicher Höhe gegenüber
(abgesehen von Vermögensänderungsbilanz und Restposten): Deutschland erwirbt Forderungen
bzw. Vermögen im Ausland. Ungleichgewichte zeigen sich deshalb nur in den
\emph{Teilbilanzen}, die Zahlungsbilanz insgesamt ist formal stets ausgeglichen.

\medskip
\textbf{c) Handelshemmnisse} \hfill [4~P.]

\emph{Tarifäre Handelshemmnisse} wirken über den \emph{Preis} der gehandelten Güter
(preispolitische Instrumente). \emph{Nichttarifäre Handelshemmnisse} wirken über
\emph{Mengen} oder \emph{administrative/technische Hürden}. [2~P.]

Zuordnung [je 0,5~P.]: (1)~Einfuhrzoll $\Rightarrow$ tarifär; (2)~Einfuhrkontingent
(Höchstmenge) $\Rightarrow$ nichttarifär; (3)~Sicherheitszertifizierung (technische Norm,
Verwaltungsverfahren) $\Rightarrow$ nichttarifär; (4)~Exportsubvention $\Rightarrow$
tarifär.

\medskip
\textbf{d) Beurteilung von Gegenzöllen} \hfill [4~P.: je Argument 1~P.]

\emph{Dafür} (z.\,B.):
\begin{itemize}
  \item Druckmittel, um Westland an den Verhandlungstisch zu bringen und die Zölle
    zurückzunehmen.
  \item Schutz europäischer Hersteller vor subventionierten Billigimporten aus
    Westland (Herstellung fairer Wettbewerbsbedingungen).
\end{itemize}
\emph{Dagegen} (z.\,B.):
\begin{itemize}
  \item Gefahr einer Eskalation (Handelskrieg) mit weiteren Zöllen auf EU-Waren ---
    gerade für ein exportorientiertes Unternehmen wie Weinmann ein großes Risiko.
  \item Importe aus Westland (Vorprodukte, Konsumgüter) verteuern sich; Kosten für
    Unternehmen und Preise für Verbraucher steigen, Wohlfahrtsgewinne der
    internationalen Arbeitsteilung gehen verloren.
\end{itemize}
Hinweis: Über Zölle entscheidet nicht Deutschland, sondern die \textbf{EU} (gemeinsame
Handelspolitik, gemeinsamer Außenzoll); Maßnahmen müssen mit den WTO-Regeln vereinbar
sein. Ein begründetes eigenes Fazit wird erwartet.
\end{loesung}

\begin{loesung}[title={Lösung Aufgabe~3.8 --- Emissionshandel im Unternehmen}]
\textbf{a) Funktionsweise und Prinzip} \hfill [3~P.]

Der Staat bzw. die EU legt eine \textbf{Obergrenze} für die Gesamtemissionen der
erfassten Anlagen fest (Cap), die schrittweise sinkt. In dieser Höhe werden
Emissionsberechtigungen (Zertifikate) ausgegeben bzw. versteigert. Ein Unternehmen darf nur
so viel CO\textsubscript{2} ausstoßen, wie es Berechtigungen abgeben kann. Fehlende
Berechtigungen müssen am Markt gekauft werden, nicht benötigte können verkauft werden
(Trade). So entsteht ein Marktpreis für CO\textsubscript{2}, der zur Vermeidung dort
anreizt, wo sie am günstigsten ist. [2~P.]

Prinzip: \textbf{Verursacherprinzip}; externe Kosten werden internalisiert. [1~P.]

\medskip
\textbf{b) Kostenvergleich} \hfill [5~P.]

\emph{Ohne Anlage:} Fehlmenge $46.000 - 40.000 = 6.000$ Berechtigungen;\\
Zukaufkosten $6.000 \times 80~\text{€} = \textbf{480.000~€}$ [1,5~P.]

\emph{Mit Anlage:} Emissionen $46.000 - 9.000 = 37.000$~t; Überschuss
$40.000 - 37.000 = 3.000$ Berechtigungen;\\
Verkaufserlös $3.000 \times 80~\text{€} = 240.000~\text{€}$;\\
Nettobelastung $540.000 - 240.000 = \textbf{300.000~€}$ [2~P.]

\emph{Entscheidung:} Die Wärmerückgewinnungsanlage ist um $480.000 - 300.000 =
\textbf{180.000~€}$ pro Jahr günstiger und sollte realisiert werden. [1~P.]\\
Kontrolle über die Vermeidungskosten: $540.000 \div 9.000 = 60$~€/t $<$ 80~€/t
Zertifikatspreis; Vorteil $9.000 \times (80 - 60) = 180.000$~€. [0,5~P.]

\medskip
\textbf{c) Beurteilung Emissionshandel vs. Auflage} \hfill [3~P.: je Argument 1~P.]

\emph{Vorteile des Emissionshandels} (zwei genügen):
\begin{itemize}
  \item \emph{Ökologische Treffsicherheit:} Die Gesamtmenge der Emissionen ist durch das
    Cap sicher begrenzt.
  \item \emph{Kosteneffizienz:} Vermieden wird dort, wo es am billigsten ist (wie bei
    Hagedorn mit 60~€/t); das Umweltziel wird volkswirtschaftlich zu geringsten Kosten
    erreicht.
  \item \emph{Dauerhafter Innovationsanreiz:} Jede eingesparte Tonne bringt Geld; bei
    einer Auflage fehlt der Anreiz, den Grenzwert zu unterschreiten.
\end{itemize}
\emph{Nachteil} (einer genügt): schwankende Zertifikatspreise erschweren die Planung;
Gefahr der Abwanderung in Länder ohne CO\textsubscript{2}-Preis (Carbon Leakage);
hoher Verwaltungs- und Kontrollaufwand (Messung, Berichtspflichten). Eine Auflage wirkt
dagegen direkt und sicher beim einzelnen Werk, ist aber unflexibel.
\end{loesung}
